\section{Variance risk and integrated variance: when certificates stay wide}\label{sec:variance_risk}

The sub-tick VIX bounds of Sections~\ref{sec:vix_futures}--\ref{sec:vix_options} raise a natural question:
does the same framework produce equally tight bounds for realised-variance derivatives under rough
volatility? Examining two spot-variance functionals---integrated variance in the tractable rough
Stein--Stein model and integrated variance in rough Bergomi---demonstrates that the answer is no, and
shows how the exact residual share $P_\mu/\Var(\Lambda_\mu)$ diagnoses the contrast before any simulation.

\subsection{Rough Stein--Stein variance risk against an exact characteristic-function law}\label{sec:case_steinstein}
Here the conditional-variance principle behind \RoughCertCV{} supplies statistical risk brackets. Their sampling uncertainty is distinct from the deterministic Arb pricing enclosures. In the rough Stein--Stein model, spot volatility is Gaussian, $\sigma_t=\sigma_0+\nu W_t^H$ with $W^H$
Riemann--Liouville fractional Brownian motion, and the characteristic function of integrated variance
$\Lambda=\int_0^T\sigma_t^2\dd t$ is known in closed form \citep{AbiJaber2022}. Because $\Lambda$ is a
quadratic subordinated functional, both the conditional-expectation monitor $\hat\Lambda=\E[\Lambda\mid\mathcal{G}]$
and the conditional variance $V=\Var(\Lambda\mid\mathcal{G})$ given $n$ monitoring dates are explicit
quadratic forms in the whitened observations,
$\hat\Lambda=d_0+\ell^\top\xi+\xi^\top A\xi$ and $V=c_0+b^\top\xi+\xi^\top\Omega\xi$, computed once by
graded quadrature. We test one month of daily closes ($T=1/12$, $n=21$, $\sigma_0=0.2$, $\nu=0.1$,
$H\in\{0.05,0.10,0.30\}$) on the $99\%$ quantile (value at risk of a short variance swap) and calls
$\E(\Lambda-k)_+$ at $k=\E\Lambda$ and $k=\E\Lambda+2\,\mathrm{sd}$ ($50$ replications of $K=2\times10^5$
daily paths, DKW level $a=0.01$, \lab{statistical}$(0.99)$). An independent Karhunen--Lo\`eve expansion
($2{,}500$ nodes, Fourier inversion) provides the reference law, matching the exact mean to $10^{-9}$ and
the ledger variance to within \EthirteenKLvar\% (omitted tail share $5.8\times10^{-6}$ at $H=0.05$).

Table~\ref{tab:e13} shows that daily monitoring leaves substantial continuum bias: the exact RMS error of
the daily Riemann sum is \EthirteenRiLo--\EthirteenRiHi\% of $\mathrm{sd}(\Lambda)$, biasing standard
daily Monte Carlo by \EthirteenNaiveLo\ to \EthirteenNaiveHi\ standard deviations on the $99\%$ quantile
and by \EthirteenPriceBiasLo\% to \EthirteenPriceBiasHi\% on calls. Every certified bracket attained
empirical coverage \EthirteenCovMin. Conditioning on daily closes with conditional variance and exact
binomial order statistics (Proposition~\ref{prop:condvar}(ii)) shrank the $99\%$ VaR brackets to
\EthirteendLo--\EthirteendHi\ standard deviations (sampling floor \EthirteendFloorHi), and computing
$\E(\hat\Lambda-k)_+$ from its quadratic-form characteristic function shrank ATM call brackets to
\EthirteencAtmLo--\EthirteencAtmHi\% of the exact price.

For $H\le0.10$, daily closes alone still leave $\ResShareDailyLo$--$\ResShareDailyHi\%$ of $\Var(\Lambda)$
unexplained. Simulating
intraday volatility paths conditionally on the daily closes via an exact Gaussian bridge leaves the
continuum law of $\Lambda$ unchanged while refining $\mathcal{G}$ (Table~\ref{tab:e13}, Panel~C): moving from one
to \EthirteenbMaxDaya\ points per day reduces the residual share from \EthirteenbSaOne\% to
\EthirteenbSaEight\% ($H=0.05$) and from \EthirteenbSbOne\% to \EthirteenbSbEight\% ($H=0.10$), cutting
the estimated conditional Cantelli $99\%$ bracket width from approximately \EthirteenbDaOne\ to \EthirteenbDaEight\ standard
deviations ($H=0.05$) and from \EthirteenbDbOne\ to \EthirteenbDbEight\ ($H=0.10$) while preserving
$1.00$ coverage. Thus the ledger quantifies the reduction in the residual variance from each tested monitoring refinement before
pricing a portfolio.

\begin{table}[!htbp]\centering\scriptsize\setlength{\tabcolsep}{3.5pt}
\caption{Rough Stein--Stein one-month integrated variance ($n=21$ daily closes): statistical $99\%$ VaR brackets (in exact sd), variance-call widths (relative to exact price), and intraday Gaussian-bridge refinement.}\label{tab:e13}
\begin{tabular}{@{}rrrcccccr@{}}\toprule
\multicolumn{9}{@{}l}{\emph{Panel A: Daily-close $99\%$ VaR bracket widths in exact sd (empirical coverage in parentheses)}} \\
\addlinespace[1.5pt]
 & \multicolumn{2}{c}{RMS error (\% of sd)} & \multicolumn{2}{c}{Baseline width (sd)} & \multicolumn{3}{c}{\textbf{\RoughCertCV{}} and sampling (sd)} & \\
\cmidrule(lr){2-3}\cmidrule(lr){4-5}\cmidrule(lr){6-8}
$H$ & Daily sum & Monitor & Riemann $W_2$ & Martingale & Cantelli & DKW floor & \shortstack{Binomial width\\and floor} & \shortstack{Daily bias\\(sd)} \\\midrule
0.05 & 42.7\% & 37.0\% & 8.50 (1.00) & 5.39 (1.00) & \textbf{\boldmath 3.45} (1.00) & 0.38 & \textbf{\boldmath 2.61} (1.00)/0.06 & $+$0.315 \\
0.10 & 25.0\% & 23.0\% & 5.70 (1.00) & 3.72 (1.00) & \textbf{\boldmath 1.99} (1.00) & 0.41 & \textbf{\boldmath 1.45} (1.00)/0.07 & $+$0.165 \\
0.30 & 7.7\% & 6.5\% & 2.54 (1.00) & 1.59 (1.00) & \textbf{\boldmath 0.71} (1.00) & 0.39 & \textbf{\boldmath 0.37} (1.00)/0.06 & $+$0.093 \\
\midrule
\multicolumn{9}{@{}l}{\emph{Panel B: Variance calls $\E(\Lambda-k)_+$ --- certified width / exact price (coverage in parentheses)}} \\
\addlinespace[1.5pt]
$H$ & Strike $k$ & Exact price & Lipschitz & Convex & Scarf--Lo & \shortstack{Exact\\monitor law} & & \shortstack{Daily bias\\(\%)} \\\midrule
0.05 & $\E\Lambda$ & $4.7\times 10^{-4}$ & 2.74 (1.00) & 1.04 (1.00) & \textbf{\boldmath 0.751} (1.00) & \textbf{\boldmath 0.177} (1.00) & & $+$11.1\% \\
0.05 & $+2\,\mathrm{sd}$ & $2.3\times 10^{-5}$ & 55.10 (1.00) & 20.67 (1.00) & \textbf{\boldmath 12.923} (1.00) & \textbf{\boldmath 1.702} (1.00) & & $+$61.0\% \\
0.10 & $\E\Lambda$ & $5.2\times 10^{-4}$ & 1.78 (1.00) & 0.80 (1.00) & \textbf{\boldmath 0.589} (1.00) & \textbf{\boldmath 0.083} (1.00) & & $+$6.0\% \\
0.10 & $+2\,\mathrm{sd}$ & $2.8\times 10^{-5}$ & 32.97 (1.00) & 14.60 (1.00) & \textbf{\boldmath 9.900} (1.00) & \textbf{\boldmath 0.756} (1.00) & & $+$28.6\% \\
0.30 & $\E\Lambda$ & $3.8\times 10^{-4}$ & 1.01 (1.00) & 0.71 (1.00) & \textbf{\boldmath 0.632} (1.00) & \textbf{\boldmath 0.011} (1.00) & & $+$3.7\% \\
0.30 & $+2\,\mathrm{sd}$ & $1.7\times 10^{-5}$ & 21.82 (1.00) & 15.04 (1.00) & \textbf{\boldmath 13.366} (1.00) & \textbf{\boldmath 0.123} (1.00) & & $+$17.5\% \\
\midrule
\multicolumn{9}{@{}l}{\emph{Panel C: Intraday Gaussian-bridge conditioning refinement for $99\%$ VaR (widths in continuum sd)}} \\
\addlinespace[1.5pt]
$H$ & \shortstack{Observations\\per day} & \shortstack{Residual\\share} & & Martingale & Cantelli & DKW floor & & \shortstack{Cantelli\\coverage} \\\midrule
0.05 & 1 & $1.4\times 10^{-1}$ & & 5.39 & \textbf{\boldmath 3.46} & 0.38 & & 1.00 \\
0.05 & 2 & $7.1\times 10^{-2}$ & & 4.12 & \textbf{\boldmath 2.30} & 0.39 & & 1.00 \\
0.05 & 4 & $3.5\times 10^{-2}$ & & 3.16 & \textbf{\boldmath 1.59} & 0.41 & & 1.00 \\
0.05 & 8 & $1.7\times 10^{-2}$ & & 2.46 & \textbf{\boldmath 1.15} & 0.40 & & 1.00 \\
0.10 & 1 & $5.3\times 10^{-2}$ & & 3.72 & \textbf{\boldmath 1.99} & 0.42 & & 1.00 \\
0.10 & 2 & $2.4\times 10^{-2}$ & & 2.82 & \textbf{\boldmath 1.38} & 0.42 & & 1.00 \\
0.10 & 4 & $1.1\times 10^{-2}$ & & 2.18 & \textbf{\boldmath 1.02} & 0.42 & & 1.00 \\
0.10 & 8 & $4.8\times 10^{-3}$ & & 1.73 & \textbf{\boldmath 0.80} & 0.43 & & 1.00 \\
\bottomrule\end{tabular}
\smallskip\par\raggedright\scriptsize\emph{Notes:} Bold identifies conditional-variance widths, rounded up. Parentheses give empirical coverage. The monitor is conditional expectation; sd is the continuum standard deviation. DKW denotes the Dvoretzky--Kiefer--Wolfowitz sampling floor. The binomial cell gives width (coverage) / floor. Daily bias compares the daily estimate with the reference law. Panels A/B use 50 replications and C uses 10.
\end{table}
%

\subsection{Rough Bergomi integrated variance: the structural contrast with VIX}\label{sec:case_iv}
For one-month integrated variance $\mathrm{IV}=\int_0^{1/12} v_t\dd t$ and integrated volatility
$\mathrm{IVol}=\int_0^{1/12}\sqrt{v_t}\dd t$ under rough Bergomi ($\xi_0=0.235^2$, $\eta_{\mathrm{P1}}=1.9$,
$H\in\{0.05,0.10,0.30\}$), the spot variance path is integrated directly rather than smoothed by a forward
expectation. Conditioning on the $n=21$ daily closes requires a singularity-corrected per-sample rule for
$V(\xi)=\iint\hat X_s\hat X_u(e^{\eta^2 D(s,u)}-1)\dd s\dd u$ (maximum relative per-sample audit difference
\EsixcAudit, added conservatively to $v$; \lab{statistical}$(0.99)$ over $K=2\times10^5$ paths).

Table~\ref{tab:e6c} summarises the outcome. Although conditional variance narrows the $99\%$ quantile
brackets by factors \EsixcRatio\ over martingale-sharp bounds and the conditional Scarf--Lo bound narrows
variance-call brackets by factors \EsixcPriceGainLo--\EsixcPriceGainHi\ over convex-order bounds (with
fine-grid Monte Carlo diagnostics lying inside \EsixcInside/\EsixcCells\ quantile brackets and
\EsixcPriceIn/\EsixcPriceN\ call brackets), the absolute widths remain $\EsixcWidthLo$--$\EsixcWidthHi$ standard deviations
on daily-close quantiles. The separate 256-date quantile widths appear in Table~\ref{tab:e6c}, Panel~A. These brackets serve as model
validation envelopes, not quoting intervals.

The exact residual share explains why rough Bergomi produces both the narrowest and the widest brackets in
this study. On the VIX forward window $[T,T+\Delta]$, every forward variance $\xi_T(u)$ averages the same
past Brownian history on $[0,T]$ through a kernel smooth across the diagonal, so a single Gaussian
direction leaves at most $\ShareVIXBase\%$ of $\Var(\mathrm{VIX}_T^2)$ in base regimes ($\ShareVIXStress\%$
in six-month $\eta_{\mathrm{P1}}=3$ stress). On spot integrated variance $\int_0^T v_t\dd t$, by contrast, fresh rough
innovations arrive at every instant $t\in[0,T]$, so even $21$ daily closes leave up to $\ShareIVHi\%$ of
$\Var(\mathrm{IV})$ unexplained at $H=0.05$. Roughness alone therefore does not determine certifiability;
the computable residual share $P_\mu/\Var(\Lambda_\mu)$ separates tractable low-rank functionals from
path-dominated ones before any simulation is launched.

\begin{table}[!htbp]\centering\footnotesize\setlength{\tabcolsep}{2.7pt}
\caption{Rough Bergomi one-month integrated variance and volatility from $21$ daily closes: statistical $99\%$ quantile widths (in sd) and median variance-call bounds ($\times 10^4$).}\label{tab:e6c}
\begin{tabular}{@{}llrcccccc@{}}\toprule
\multicolumn{9}{@{}l}{\emph{Panel A: $99\%$ quantile widths in exact sd ($21$ daily closes vs.\ $256$-date refinement)}} \\
\addlinespace[1.5pt]
 & & & \multicolumn{4}{c}{21-date quantile widths (sd) and ratio} & \multicolumn{2}{c}{Diagnostics} \\
\cmidrule(lr){4-7}\cmidrule(lr){8-9}
Functional & $H$ & \shortstack{Monitor RMS\\(\% of sd)} & Martingale & \textbf{Cantelli} & DKW floor & \shortstack{Width\\ratio} & \shortstack{256 dates\\(sd)} & \shortstack{Monte Carlo\\in band} \\\midrule
Integrated variance & 0.05 & 44\% & 7.38 & \textbf{\boldmath 6.00} & 0.88 & 1.23$\times$ & 3.96 & 1.00 \\
Integrated volatility & 0.05 & 38\% & 5.91 & \textbf{\boldmath 4.16} & 0.54 & 1.42$\times$ & 2.65 & 1.00 \\
Integrated variance & 0.10 & 30\% & 6.12 & \textbf{\boldmath 4.60} & 1.12 & 1.33$\times$ & 3.16 & 1.00 \\
Integrated volatility & 0.10 & 24\% & 4.36 & \textbf{\boldmath 2.72} & 0.62 & 1.60$\times$ & 1.99 & 1.00 \\
Integrated variance & 0.30 & 7\% & 2.64 & \textbf{\boldmath 1.65} & 0.96 & 1.60$\times$ & 1.42 & 1.00 \\
Integrated volatility & 0.30 & 7\% & 1.93 & \textbf{\boldmath 0.98} & 0.56 & 1.98$\times$ & 0.89 & 1.00 \\
\midrule
\multicolumn{9}{@{}l}{\emph{Panel B: Variance-call brackets ($\times 10^4$) vs.\ fine-grid Monte Carlo diagnostic}} \\
\addlinespace[1.5pt]
$H$ & Strike & & \multicolumn{2}{c}{Convex order} & \multicolumn{2}{c}{\textbf{Conditional Scarf--Lo}} & \shortstack{Fine-grid\\estimate} & \\\midrule
0.05 & $\E\,\mathrm{IV}$ &  & \multicolumn{2}{c}{[10.362,\,22.828]} & \multicolumn{2}{c}{\bfseries\boldmath [10.362,\,17.512]} & 12.265 &  \\
0.05 & $+2\,\mathrm{sd}$ &  & \multicolumn{2}{c}{[0.859,\,13.037]} & \multicolumn{2}{c}{\bfseries\boldmath [0.859,\,6.729]} & 2.064 &  \\
0.10 & $\E\,\mathrm{IV}$ &  & \multicolumn{2}{c}{[11.898,\,23.209]} & \multicolumn{2}{c}{\bfseries\boldmath [11.898,\,18.427]} & 13.552 &  \\
0.10 & $+2\,\mathrm{sd}$ &  & \multicolumn{2}{c}{[1.319,\,12.335]} & \multicolumn{2}{c}{\bfseries\boldmath [1.319,\,7.068]} & 2.494 &  \\
0.30 & $\E\,\mathrm{IV}$ &  & \multicolumn{2}{c}{[9.107,\,14.147]} & \multicolumn{2}{c}{\bfseries\boldmath [9.107,\,13.256]} & 9.824 &  \\
0.30 & $+2\,\mathrm{sd}$ &  & \multicolumn{2}{c}{[0.925,\,5.755]} & \multicolumn{2}{c}{\bfseries\boldmath [0.925,\,4.822]} & 1.415 &  \\
\bottomrule\end{tabular}
\smallskip\par\raggedright\scriptsize\emph{Notes:} Bold identifies \RoughCertCV{} conditional Cantelli or Scarf--Lo results. Monitor RMS is root-mean-square error relative to the continuum standard deviation. The width ratio divides martingale by Cantelli width. DKW is the Dvoretzky--Kiefer--Wolfowitz sampling floor. Monte Carlo in band is the fraction of five replications enclosing the fine-grid diagnostic ($200{,}000$ samples per replication).
\end{table}
%
\ifTwoColumn\else\FloatBarrier\fi
